% Generated by scripts/assemble.py; corrections belong in source/corrections.json.
\examhead{34}{复旦大学 2025 年数学分析部分}{37}
% p038-b17; source PDF page 38
\par\Needspace{5\baselineskip}
\blockanchor{p038-b17}
\majorquestion{一}填空题. 每题 10 分, 共 50 分.\par
% p038-b18; source PDF page 38
\blockanchor{p038-b18}
\question{1}求极限  $\lim_{n\to\infty}\left(\frac{1+\sqrt[n]{2025}}{2}\right)^n=$  \_\_\_\_.\par
% p038-b19; source PDF page 38
\blockanchor{p038-b19}
\question{2}求不定积分  $\int \frac{1 + \sin x}{1 + \cos x} e^{x} \, dx =$ \_\_\_\_ .\par
% p038-b20; source PDF page 38
\par\Needspace{6\baselineskip}
\blockanchor{p038-b20}
\question{3}$f(x)$ 为定义在 $(0,1)$ 上的实值函数，满足\par
\nopagebreak[4]
% p038-b21; source PDF page 38
\blockanchor{p038-b21}
\displayformula{\lim _ {x \rightarrow 0 ^ {+}} f (x) = 0, \lim _ {x \rightarrow 0 ^ {+}} \frac {f (x) - f \left(\frac {x}{2}\right)}{x} = 0.}
% p038-b22; source PDF page 38
\blockanchor{p038-b22}
则  $\lim_{x\to0^{+}}\frac{f(x)}{x}=$  \_\_\_\_.\par
% p038-b23; source PDF page 38
\blockanchor{p038-b23}
\question{4}平面 $x + y + z = 0$ 与椭球面 $x^{2} + y^{2} + 4z^{2} = 1$ 相交得到的椭圆面积 $S =$ \_\_\_\_.\par
% p039-b00; source PDF page 39
\blockanchor{p039-b00}
\question{5}设 L 为任一包含原点的光滑曲线, 取逆时针方向, 求  $\oint_{L}\frac{x\,dy-y\,dx}{x^{2}+y^{2}}=$  \_\_\_\_.\par
% p039-b01; source PDF page 39
\par\Needspace{5\baselineskip}
\blockanchor{p039-b01}
\majorquestion{二}(20 分) 设 $\alpha > \beta > 1$, 求 $\int_{0}^{+\infty} \frac{1}{1 + x^{\alpha} |\sin x|^{\beta}} \, \mathrm{d}x$ 的敛散性.\par
% p039-b02; source PDF page 39
\par\Needspace{5\baselineskip}
\blockanchor{p039-b02}
\majorquestion{三}(20 分) 求使级数 $\sum_{n=1}^{\infty} \left( \frac{1}{n} \csc \frac{1}{n} - 1 \right)^{2025p}$ 收敛的 $p$ 的取值范围.\par
